\chapter{93}\label{section-93}

Louis schloss die Tür zum Büroraum und schlich sich in die Küche hinter
die Eingangstür. Die Vier sassen offensichtlich um den Tisch vor dem
Haus und waren in ein lebhaftes Gespräch vertieft. Bevor er hinaustrat,
wollte er wissen, worüber sie sprachen. Bens Stimme hörte er zuerst.

«Das dachte ich auch, Manuela. Nur, die Gewalt durch Frauen oder Mädchen
macht nicht mehr Halt vor unseren Türen. Natürlich zeigen Statistiken
diesen Trend schon einige Jahre mit erschütternder Klarheit. Nur, wir
studierten Mädchengangs bisher aus Distanz im nahen Ausland, bis eben
nun auch bei uns ein Fall dieser Tragweite eintrat. Und dann wurde
tatsächlich \emph{ich} als anwaltlicher Verteidiger für die Gang
angefragt. Der Fall dauerte fast zwei Jahre. Ich ahnte damals, was auf
mich zukommen würde, trotzdem...»

Joh unterbrach ihn.

«...hast du ihn angenommen, B.B.»

«Ja, aber nur, weil mir Claudine beiläufig mitteilte, dass sie das
geschädigte Mädchen vertritt. Ich sah eine gute Chance der
Zusammenarbeit. Nur, ich sag euch,» Ben seufzte wie ein müder, alter
Mann, «vorläufig will ich nun einfach mal leben, entspannt und zusammen
mit Franziska schöne Dinge tun und glücklich sein.»

«Bingo, endlich, Dad, Mam wird sich freuen.»

Ava drückte ihrem Vater einen Kuss auf die Wange.

«Hat sie dir davon...»

«Nein, Dad, hat sie nicht,» Ava verdrehte die Augen, nahm seine Hand und
drückte sie, «aber etwas möchte ich echt wissen. Wie erklärst du dir den
Ursprung dieser abscheulichen Behandlung von diesem Mädchen durch
Gleichaltrige? Mit dreizehn, vierzehn sind das doch noch Kinder, und als
Mädchen, pff,» Ava schüttelte den Kopf, «wie kommt das?»

«Ganz abgesehen von den sozialen Hintergründen der Angreiferinnen spielt
das gleiche Phänomen wie bei männlichen Tätern. Es ist ein eindeutiger
Fall von \emph{Verantwortungsdiffusion.»}

Ava zog die Brauen hoch. Manuela rückte näher an den Tisch heran und
legte die Hände übereinander.

«Und das heisst?»

«Ganz einfach erklärt, die Gewalt geht von der Gruppe aus, das heisst,
die Wutdynamik rechtfertigt die eigene Gewalt. Gewalt steckt an und vor
allem, sie wird legitim, die anderen tun’s ja auch.»

Joh strich sich nachdenklich über die Oberlippe.

«So funktionieren Kriege.»

Ben nickte.

Louis fasste Mut. Jetzt musste er raus. Vielleicht würden sie
weiterdiskutieren und sein Dazustossen könnte nebenher geschehen. Als er
einen Schritt nach vorne gehen wollte, begann es an seinem Bein zu
vibrieren. Louis trat zurück und zog das Handy aus der Seitentasche
seiner Hose. Wieder ein Anruf mit unterdrückter Nummer. Kalter Schweiss
setzte sich auf seiner Stirn fest. Seine Beine wurden schwer und ein
Zittern regierte beide Hände. Im gleichen Moment hörte er Schritte die
Treppe hochkommen und Joh stand vor ihm.

«Ciro, ach, hier bist du, komm doch auch,» Joh hielt abrupt inne und
starrte Louis entgeistert an, «was ist denn mit dir los?»

Louis konnte sich kaum mehr auf den Beinen halten. Um ein Haar wäre er
Joh entgegengefallen.

«Komm,» Joh stützte ihn und führte ihn langsam zur Sofaecke des
Wohnraumes, «setz dich, Ciro, ich bin gleich bei dir. Ich gebe draussen
kurz Bescheid, dass sie nicht auf mich warten sollen.»

Louis wusste es. Es war soweit. Die Zeit der Wahrheit war jetzt.

Joh war rasch zurück. Er setzte sich neben Louis. Dann sassen sie da.
Nebeneinander. Joh blieb still und wartete.

Eine unergründliche Ruhe breitete sich in Louis aus. Die Luft schien
fein zu flirren. Er nahm nur noch ein Gefühl von Wärme wahr. Sein Atem
beruhigte sich. Von Joh ging eine Kraft aus, die Louis in sich
zusammenfallen liess. Er verstand nicht, wie es kam, dass er sich mehr
und mehr entspannte. Gerade jetzt, wo er Joh gewaltig enttäuschen würde.

Es mochte ewig gedauert haben, bis Louis es wagte, hoch und ganz langsam
zu Joh hinüber zu blicken. So lange und direkt hatte er bisher nur
wenige Menschen angesehen. Johs Gesicht zeigte eine Gutmütigkeit, die
ihm fremd geworden war. Es war unmöglich, wegzusehen. Louis’ Augen
wurden feucht. Schliesslich flossen Tränen über seine Wangen. Und Joh
schaute unbeweglich zu. So lange, bis Louis’ Körper zu vibrieren begann
und ihn ein Zittern einnahm, das er nicht kontrollieren konnte. Dann
öffnete Joh seine Arme. Seine Stimme klang weich.

«Komm, Ciro.»

Mit den Worten zog er Louis näher an sich heran und schloss ihn in die
Arme.

Bruchstückweise, ähnlich auseinandergeschnittener Sätze, die es neu
zusammenzusetzen galt, erzählte Louis seine Geschichte. Einige Male
blickte er zu Joh hoch, legte aber seinen Kopf gleich wieder an dessen
Brust zurück. Er hatte es nur knapp geschafft, Joh die Stockszene zu
verschweigen.

Es war dieses geborgene Gefühl, wie es Louis vor ein paar Tagen
empfunden hatte. Diesmal aber wäre er in Johs Umarmung am liebsten
eingeschlafen und nicht mehr aufgewacht.

\vspace{\baselineskip}
\emph{Damals, Fribourg, 10

Vater bringt Louis zu Maman zurück. Nach nur einer halben Stunde. Wegen
dieses blöden Anrufs. «Der Kunde ist König,» raunt ihm der Vater zu,
stellt ihn wieder vor der Haustür ab, wo er ihn doch eben erst abgeholt
hat und während er die Autotür öffnet, ruft er ihm zu: «Ich melde mich
bald wieder, mein Lieblingssohn.»

Es geht alles so schnell, dass Louis nicht weiss, wie ihm geschieht. Er
steht wieder vor der Haustür, spürt die Tränen kommen und mit ihnen eine
richtig grosse Wut. Die Wohnungstür schmettert er hinter sich ins
Schloss, schiebt seine erstaunte Schwester grob zur Seite, flüchtet in
sein Zimmer und dreht den Schlüssel. Er wirft sich auf sein Bett und
drückt das Gesicht ins Kopfkissen. Dann hört er ein leises Fiepsen.
Teddy. Nicht auch noch. Dieses blöde Viech.

Louis schnellt hoch, öffnet das Türchen des Käfigs, nimmt den Hamster
ruckartig raus und auf die Hand. Dann klopft es an seine Zimmertür.

«Louis?»

Es ist Maman.

«Lass mich.»

«Louis, bitte mach die Tür auf.»

«Lass mich.»

Nach ein paar Sekunden hört er Mamans Schritte, die leiser werden.

Der Hamster fiepst weiter. Dann drückt Louis zu. Immer stärker. Das Tier
wehrt sich. Es fällt ihm aus der Hand, knallt entlang des Stuhlbeins
nach unten und landet auf dem Plattenboden. Louis starrt es an. Es
bewegt sich nicht mehr. Er geht in die Knie und glaubt einzufrieren.

«Louis, bitte lass mich rein. Ich versteh doch, dass du traurig bist.»

Es ist erneut Mamans Stimme.

Louis hebt den Hamster hoch. Jetzt bewegt der den Kopf. Louis atmet auf,
legt ihn ganz schnell zurück in den Käfig und schliesst das Türchen.
Dann schiebt er den Käfig unter sein Pult und legt die Tagesdecke
darüber. Er schämt sich. Maman und Lucia dürfen nichts wissen. Er hofft,
dass Teddy nichts Schlimmes passiert ist.

Zwei Tage geht alles gut. Louis vermeidet es, den Käfig zu reinigen. Die
Fütterung hält er kurz und beim Schliessen des Türchens schafft er es,
den Hamster zu ignorieren. Teddy frisst zwar wenig, bewegt sich
auffallend langsam, aber dennoch, er lebt.

Bis es Louis sieht, an diesem einen Abend, ehe er ins Bett steigt. Das
linke Hinterbein des Tieres steht ab. Es liegt unnatürlich abgezweigt
neben ihm, als gehöre es nicht zu Teddy. Louis reisst die Augen auf. Er
wagt nicht, Teddy anzufassen. Dann hält er die Luft an. Ein beissend
schlechter Geruch steigt aus dem Käfig zu ihm hoch. Louis hört Mamans
Schritte. Sie wird zu ihm kommen, wie jeden Abend, und sich an sein Bett
setzen. Sie werden über den Tag reden und Maman wird ihm eine gute Nacht
wünschen. Was nun? Louis wird panisch. Soll er die Decke einfach über
den Käfig legen?

«So, Louis, da bin ich.»

Maman setzt sich an seinen Bettrand und lächelt. Plötzlich kneift sie
die Augen zusammen und schaut sich suchend um.

«Du, sag mal, was riecht denn hier so streng?»

Jetzt kann Louis nicht mehr. Er wirft sich Maman entgegen, hängt sich an
sie und weint. Sie hält ihn.

«Wenn er stirbt, dann, Maman, dann habe ich, hab ich ihn ermordet.»

Louis drückt sich ganz fest an sie.

«Wen hast du ermordet, Louis, was redest du da?»

«Teddy, Maman, Teddys Bein ist fast ab. Ich war, ich war so wütend.»

Maman hört lange zu. Bis die ganze Geschichte raus ist. Sie schimpft
kein einziges Mal. Dann löst sich Louis aus der Umarmung und wischt sich
die Tränen mit dem Ärmel trocken.

Teddy lebt noch knappe zwölf Tage.}
\vspace{\baselineskip}

Zusammen mit Maman und Lucia gräbt Louis ein ganz tiefes Loch im Wald.
Er legt Teddy hinein und bedeckt ihn mit Erde, bis der Waldboden wieder
geschlossen ist.
